\scanpage{126}
\mspage{-III- \handwritten{7} -}
\begin{flushright}number \handwritten{383}\end{flushright}
quotient an essential group in the Euclidean space $E/W'$.

We say that $G$ is \uline{reducible} if $E$ is the direct sum of two orthogonal subspaces $E',E''$ distinct from $E$ and stable under ${}^vG$. Then every element $s$ of $G$ defines, by passage to the quotient, affine automorphisms $s'$ of $E/E''$ (identified with $E'$) and $s''$ of $E/E'$ (identified with $E''$); the set of the $s'$ (resp. $s''$) is a subgroup $G'$ (resp. $G''$) of the affine group of $E'$ (resp. $E''$). If $H$ is a hyperplane of $E$ such that $s_H\in G$, one easily sees that either $E'$ is contained in the direction ${}^vH$ of $H$, or $E'$ contains the line orthogonal to $H$; in the first case the automorphism $s'$ of $E'$ defined by $s_H$ is the identity, and in the second $s'$ is the reflection with respect to the hyperplane $H\handwritten{\cap}E'$ of $E'$. It follows that $G'$ and $G''$ are generated by reflections. On the other hand, the map $s\handwritten{\rightsquigarrow}(s',s'')$ from $G$ to $G'\handwritten{\times}G''$ is obviously an isomorphism.

We say that $G$ is \uline{irreducible} if it is not reducible. Denoting by $\underline{H}$ the set of hyperplanes such that $s_H\in G$, this amounts to saying that there is no partition of $\underline{H}$ into two subsets such that every hyperplane of the first is perpendicular to every hyperplane of the second. In the general case it follows from what we have just seen that $E$ is khanonically identified with a product of Euclidean spaces $E^{(1)},\ldots,E^{(q)}$, and $G$ with a product of groups $G^{(1)},\ldots,G^{(q)}$ acting on these spaces, in such a way that each $G^{(i)}$ is irreducible.

We shall henceforth fix a set $\underline{H}$ of (affine) hyperplanes of $E$ such that

1) the reflections $s_H$ $(H\in\underline{H})$ generate $G$,

2) the hyperplanes of $\underline{H}$ are permuted by $G$.

As in number 1, let us note that, for $g\in G$, we have $s_{g(H)}=gs_Hg^{-1}$. We shall see later that the set of the $s_H$ $(H\in\underline{H})$ is the set of all the reflections of $G$.

PROPOSITION 4.- \uline{The family of hyperplanes} $\underline{H}$ \uline{is locally finite}.

Let us argue by contradiction. Let $x\in E$ be such that every neighborhood $U$ of $x$ meets infinitely many hyperplanes of $\underline{H}$. Then $g(U)\handwritten{\cap}U\handwritten{\ne\varnothing}$ for infinitely many affine reflections
% archive sheet number: 125
